\documentclass[10pt,a4paper]{amsart}
\usepackage[margin=19mm]{geometry}
\usepackage{amsmath,amssymb,mathtools}
\usepackage[scr=rsfs,cal=euler]{mathalfa}
\usepackage{xfrac}
\usepackage[unicode,hidelinks,bookmarks=false]{hyperref}
\newcommand{\ie}{i.e.\ }
\newcommand{\cf}{cf.\ }
\newcommand{\resp}{respectively\ }
\newcommand{\Subsection}{\subsection}
\providecommand{\nobreakdash}{\nobreak\hbox{-}}
\setlength{\emergencystretch}{2em}
\multlinegap0pt
\usepackage{bm}
\usepackage{bbm}
\usepackage{dsfont}
\usepackage{xspace}
\usepackage{cancel}
\usepackage{mathtools}
\usepackage{amsthm}
\makeatletter
\@ifundefined{theorem}{\newtheorem{theorem}{Theorem}}{}
\makeatother
\title[Scalable Fixed-Point Framework for High-Dimensional Hamilton]{Scalable Fixed-Point Framework for High-Dimensional Hamilton-Jacobi Equations}
\author{Yesom Park, Stanley Osher}
\date{Source: arXiv:2602.05124v1 (CC BY 4.0)}
\begin{document}
\maketitle

\noindent\emph{Source and licence.} Scalable Fixed-Point Framework for High-Dimensional Hamilton-Jacobi Equations. Version arXiv:2602.05124v1 (CC BY 4.0), distributed under the Creative Commons Attribution 4.0 International licence, as recorded in the arXiv licence metadata for this exact version and shown on the abstract page https://arxiv.org/abs/2602.05124v1. Full rights notice: licenses/arxiv-2602.05124-CC-BY-4.0.txt.\par\smallskip

\noindent\textit{Editorial scope.} Counted unit: the Theorem whose source label is \texttt{thm:residual\_error} in the article (source0.tex lines 526--548, source label \texttt{thm:residual\_error}), delivered together with the article's own complete original proof of it (source0.tex lines 550--617, one \texttt{proof} environment of 68 lines). No mathematics is written, repaired or re-derived: statement and proof are the article's own text line for line. Required context retained uncounted: eq:hopf\_lax (lines 246--248); eq:fixed\_pts (lines 252--254); thm:fixed\_pt (lines 362--388). Every definition, display and label of those lines is retained unchanged. Omitted from this package and not redistributed: 1--245, 249--251, 255--361, 389--525, 549--549, 618--1287. Nothing inside a retained excerpt is omitted or altered: every retained body is delivered exactly as the article's lines are written, so no omission receipt of any kind accompanies this package. Citations: the retained excerpts carry no citation command at all, so no bibliography entry of the article is transcribed and no reference is redistributed. Reference adaptation: none was needed and none was made; every reference inside the delivered unit resolves inside it, so no unresolved reference or citation is produced. Printed slips kept literal: no printed word, symbol, hypothesis or number of the counted statement or of its proof was altered. Layout and class: the article's own class is \texttt{siamart250106} with packages including lipsum, amsfonts, graphicx, epstopdf, algorithmic, amsopn, subfigure, booktabs, empheq, multirow; the wrapper is the lane's pinned \texttt{amsart} editorial wrapper at 10pt on A4, which restates the theorem-like environments the retained text needs and adds the article's own macro definitions that the retained text needs (none), with extra wrapper packages (bm, bbm, dsfont, xspace, cancel, mathtools). The delivered numbering is the unit's own. Assets: no retained span contains a figure environment, a graphics-inclusion command, a PDF-page inclusion, a table or a TikZ picture; no image file is copied into this package, the package declares no asset dependency and no image file is redistributed with it. Licence: version arXiv:2602.05124v1 of this article is distributed under the Creative Commons Attribution 4.0 International licence, as recorded in the arXiv licence metadata for this exact version and shown on the abstract page https://arxiv.org/abs/2602.05124v1. A full rights notice is delivered at licenses/arxiv-2602.05124-CC-BY-4.0.txt. Characters: every retained excerpt is pure ASCII, so the delivered engine, XeLaTeX with its default Latin Modern text font, reports no missing character.
\begin{equation}\label{eq:hopf_lax}
    u(x,t) = \inf_{y \in \mathbb{R}^n} \Big\{ t\,H^\ast\Big(\frac{x-y}{t}\Big) + g(y) \Big\},
\end{equation}

\begin{equation}\label{eq:fixed_pts}
    y = x - t\, \nabla H(\nabla g(y)),
\end{equation}

\begin{theorem}[Fixed-point convergence]\label{thm:fixed_pt}
Let \( H: \mathbb{R}^n \to \mathbb{R} \) and \( g: \mathbb{R}^n \to \mathbb{R} \) be continuously differentiable functions such that:
\begin{itemize}
    \item \( H \) is \emph{strictly convex} and has \( L_H \)-Lipschitz continuous gradient:
    \[
    \|\nabla H(p_1) - \nabla H(p_2)\| \leq L_H \|p_1 - p_2\|, \quad \forall p_1, p_2 \in \mathbb{R}^n.
    \]
    \item \( g \) is \emph{convex} and has \( L_g \)-Lipschitz continuous gradient:
    \[
    \|\nabla g(y_1) - \nabla g(y_2)\| \leq L_g \|y_1 - y_2\|, \quad \forall y_1, y_2 \in \mathbb{R}^n.
    \]
\end{itemize}
Define the fixed-point operator
\[
    F(y) := x - t\, \nabla H(\nabla g(y)).
\]
Then, for any \(x \in \mathbb{R}^n\) and any time step \(t>0\) satisfying
\[
    t\, L_H\, L_g < 1,
\]
the iteration
\[
    y^{(k+1)} = F(y^{(k)}), \qquad k = 0,1,2,\dots
\]
converges to the unique global minimizer of the Hopf--Lax formula \eqref{eq:hopf_lax}.

\end{theorem}

\begin{theorem}[Error Analysis]\label{thm:residual_error}
Let \(H\) and \(g\) satisfy the assumptions of Theorem~\ref{thm:fixed_pt}, and let \(t>0\) satisfy \(L_F := t L_H L_g < 1\).  
Consider the iteration \(y^{(k+1)} = F(y^{(k)})\) with \(F(y) = x - t\,\nabla H(\nabla g(y))\),
and denote the minimizer by \(y^\ast\). Define the residual \(r^{(k)} := y^{(k+1)} - y^{(k)}\).
Then:
\begin{enumerate}
    \item (\textbf{Hopf--Lax Minimizer error})
    \[
    \|y^{(k)} - y^\ast\| \le \frac{1}{1 - L_F} \|r^{(k)}\|.
    \]
    \item (\textbf{Solution error})  
    If \(H^\ast\) is \(L_{H^\ast}\)-Lipschitz on the relevant domain, then
    \[
    |u(x,t) - u^{(k)}(x,t)| \le \frac{L_{H^\ast}/t + L_g}{1 - L_F}\, \|r^{(k)}\|,
    \]
    where \(u^{(k)}(x,t) := t\, H^\ast\!\big(\frac{x - y^{(k)}}{t}\big) + g(y^{(k)})\).
    \item (\textbf{Solution gradient error})  
    \[
    \|\nabla u(x,t) - \nabla u^{(k)}(x,t)\| 
\le L_g\| r^{(k)} \|.
    \]
\end{enumerate}
\end{theorem}

\begin{proof}
\textbf{1. Hopf--Lax minimizer error bound.} 
  By Theorem \ref{thm:fixed_pt}, the minimizer of the Hopf--Lax formula $y^\star$ is the unique fixed point of \(F\).
Since \(F\) is a contraction with constant \(L_F\), we have
\[
\|F(y) - F(y^*)\| \le L_F \|y - y^*\| \quad \forall y \in \mathbb{R}^n.
\]

By definition of the residual,
\[
r^{(k)} = y^{(k+1)} - y^{(k)} = F(y^{(k)}) - y^{(k)}.
\]

Then,
\[
\begin{aligned}
\|y^{(k)} - y^*\| 
&= \|y^{(k)} - F(y^*)\| \\
&= \|y^{(k)} - F(y^{(k)}) + F(y^{(k)}) - F(y^*)\| \\
&\le \|r^{(k)}\| + \|F(y^{(k)}) - F(y^*)\| \\
&\le \|r^{(k)}\| + L_F \|y^{(k)} - y^*\|.
\end{aligned}
\]

Rearranging gives
\[
\|y^{(k)} - y^*\| \le \frac{1}{1 - L_F} \|r^{(k)}\|.
\]

\medskip
\textbf{2. Solution error bound.}  
Let
\[
u(x,t) = t H^*\Big(\frac{x - y^*}{t}\Big) + g(y^*), \quad u^{(k)}(x,t) = t H^*\Big(\frac{x - y^{(k)}}{t}\Big) + g(y^{(k)}).
\]

Using the Lipschitz continuity of \(H^*\) and \(g\),
\[
\begin{aligned}
|u(x,t) - u^{(k)}(x,t)| 
&\le t \left|H^*\Big(\frac{x - y^*}{t}\Big) - H^*\Big(\frac{x - y^{(k)}}{t}\Big)\right| + |g(y^*) - g(y^{(k)})| \\
&\le t \, L_{H^*} \frac{\|y^{(k)} - y^*\|}{t} + L_g \|y^{(k)} - y^*\| \\
&= (L_{H^*}/t + L_g) \|y^{(k)} - y^*\| \\
&\le \frac{L_{H^*}/t + L_g}{1 - L_F} \|r^{(k)}\|.
\end{aligned}
\]

\medskip
\textbf{3. Solution gradient error bound.}  
Taking derivative of Hopf--Lax formula with respect to $x$ gives
\[\nabla u(x,t) = \nabla H^*((x - y^*)/t).
\] 
By incorporating the implicit relation \eqref{eq:fixed_pts}, it
gives
\[\nabla u(x,t) = \nabla H^*(\nabla H(\nabla g(y^\ast)))=\nabla g(y^\ast),
\] 
where the last equality comes from the fact $ \nabla H^*=\nabla H^{-1}$.
Consequently, we have
\[
\begin{aligned}
\|\nabla u - \nabla u^{(k)}\| 
&= \Big\|\nabla g(y^{(k)}) - \nabla g(y^\ast)\Big\| \\
&\le L_g\| r^{(k)} \|.
\end{aligned}
\]

This completes the proof.
\end{proof}

\end{document}
